\documentclass[11pt,a4paper]{article}
\usepackage[utf8]{inputenc}
\usepackage[T1]{fontenc}
\usepackage{geometry}
\usepackage{enumitem}
\usepackage{listings}
\usepackage{xcolor}
\geometry{margin=2.5cm}

\lstset{
  basicstyle=\ttfamily\small,
  breaklines=true,
  frame=single,
  backgroundcolor=\color{gray!10}
}

\title{Project Lab Logbook: Open Authorization Protocol}
\author{02244 Logic for Security\\
Fahmida Khan (s250116), Zeinab Ali, Mirtha Arul}
\date{March 16, 2026}

\begin{document}
\maketitle

\section*{Abstract}
This logbook documents the iterative development of an Open Authorization protocol.
The protocol evolves from a simplified baseline to identity-provider-based delegation, typed formats, channel abstraction, and guessable-password hardening.
Each weekly entry records: changes, modeling assumptions/goals, OFMC findings, and design decisions.

\section{Introduction}
The objective is to allow user $A$ to authorize service $p$ to access selected data on server $B$.
Authorization is mediated by trusted identity provider $I$.
The design must protect confidentiality and authentication while handling weak password assumptions for $pw(A,I)$.

\section{Week 2: Dolev--Yao Baseline (Author: Zeinab Ali)}
\textbf{What changed and idea.}
We started from a minimal three-party delegation model.
For this early baseline, authentication was simplified so that $A$ directly authenticates delegation context.

\textbf{Modeling considerations.}
Public keys were assumed available.
The model focused on obtaining a runnable baseline for static Dolev--Yao analysis.

\textbf{Problems observed.}
The baseline is intentionally weak and not yet realistic with respect to the final assignment assumptions.

\textbf{Discussion/decision notes.}
We accepted this simplification as a Week 2 starting point and planned to move trust to IDP in Week 3.

\section{Week 3: IDP and Lazy Intruder (Author: Fahmida Khan)}
\textbf{What changed and idea.}
We introduced the Identity Provider $I$.
$A$ now shares a weak password $pw(A,I)$ with $I$, and authentication uses hash evidence $h(pw(A,I),N_A)$.

\textbf{Modeling considerations.}
Password is used as authentication evidence, not as an encryption key.
This aligns with assignment constraints on weak passwords.

\textbf{Problems observed.}
OFMC still reports attacks in this stage for broader goals (\texttt{ATTACK\_FOUND}), so more hardening is required.

\textbf{Special task completion (lazy intruder).}
Role $p$ is executable with $p=i$.
Trace analysis confirms: (1) $i$ receives $\{A,p,F,T,N_A\}_{inv(pk(I))}$ from $A$; (2) $i$ composes the outgoing message to $B$ by forwarding the signed token.
This shows the service role can be automated or impersonated once a valid token is released.

\section{Week 4: Type-Flaw Resistance and GetPublicKey (Author: Fahmida Khan)}
\textbf{What changed and idea.}
To prevent type-flaw attacks, we introduced explicit function symbols so a Grant term cannot be confused with Photo Data.
We also implemented a GetPublicKey auxiliary protocol to model retrieval of $pk(B)$ and $pk(p)$ from trusted IdP $I$.

\textbf{Modeling considerations.}
Week 4 removes the unrealistic assumption that $A$ initially knows every public key.
Typed constructors enforce structure at symbolic level.

\textbf{Problems observed.}
Several heavily wrapped variants became non-executable in OFMC before simplification.
Final key-lookup model executes but still exposes weak-auth issues in analysis.

\textbf{GetPublicKey model.}
\begin{lstlisting}
Protocol: OpenAuth_W4_KeyLookup

Types:
  Agent A,I,X;
  Function pk,fKeyReq,fKeyResp;

Knowledge:
  A: A,I,X,pk(A),pk(I),inv(pk(A)),fKeyReq,fKeyResp;
  I: A,I,X,pk(A),pk(I),pk(X),inv(pk(I)),fKeyReq,fKeyResp;

Actions:
  A -> I: fKeyReq(A,X)
  I -> A: {fKeyResp(X,pk(X))}inv(pk(I))

Goals:
  A authenticates I on fKeyResp(X,pk(X));
\end{lstlisting}

\section{Week 5: Pseudonymous Channels (Author: Mirtha Arul)}
\textbf{What changed and idea.}
We created a TLS-style variant using pseudonymous channels (\texttt{*->*}) to replace as much direct transport cryptography as possible.

\textbf{Modeling considerations.}
Channel abstraction reduces protocol verbosity and clarifies trust assumptions, while preserving authorization logic.

\textbf{Problems observed.}
OFMC result is \texttt{ATTACK\_FOUND} on secrecy in this variant.
Conclusion: channel abstraction alone does not fix authorization logic flaws.

\section{Week 6: Guessable Password Security (Author: Mirtha Arul)}
\textbf{What changed and idea.}
We tested low-entropy password robustness in three steps:
baseline guessable model, hardened proof model, and focused guessable-only check.

\textbf{Modeling considerations.}
The password remains outside direct encryption usage; only hash-based witness uses $pw(A,I)$.
The final hardened flow uses pseudonymous channels between $A$ and $I$ with an explicit challenge nonce $NI$ before token issuance.

\textbf{Problems observed and outcomes.}
\begin{itemize}[leftmargin=1.2em]
  \item Baseline: \texttt{ATTACK\_FOUND}, goal \texttt{guesswhat}.
  \item Hardened model: still \texttt{ATTACK\_FOUND}, but on secrecy rather than guessability.
  \item Focused guessable-only model (\texttt{--numSess 1}): \texttt{NO\_ATTACK\_FOUND}.
\end{itemize}

\textbf{Cautious conclusion.}
Within bounded verification and for the dedicated guessable-password goal, no guessing attack was found.

\section{Final Hardened Protocol Snapshot}
\begin{lstlisting}
Protocol: OpenAuth_W6_GuessablePassword_Hardened

Types:
  Agent A,B,I,p;
  Number F,T,R,NI,Grant,Photos;
  Function pk,pw,h,fGrant,fData;

Knowledge:
  A: A,B,I,p,pk(I),pw(A,I),h,fGrant,fData;
  B: A,B,I,p,pk(B),pk(I),pk(p),inv(pk(B)),fGrant,fData;
  p: A,B,I,p,pk(p),pk(I),inv(pk(p)),fGrant,fData;
  I: A,B,I,p,pk(B),pk(p),pk(I),inv(pk(I)),pw(A,I),h,fGrant,fData;

Actions:
  p -> A: p,B,F,T,R
  A *->* I: A,p,B,F,T,R
  I *->* A: NI
  A *->* I: A,p,B,F,T,R,NI,h(pw(A,I),A,p,B,F,T,R,NI)
  I -> p: {fGrant(A,p,B,F,T,Grant)}inv(pk(I))
  p -> B: {fGrant(A,p,B,F,T,Grant)}inv(pk(I)),F
  B -> p: {fData(Photos)}pk(p)

Goals:
  pw(A,I) guessable secret between A,I;
  fData(Photos) secret between B,p;
  p authenticates I on fGrant(A,p,B,F,T,Grant);
  B authenticates I on fGrant(A,p,B,F,T,Grant);
  I authenticates A on NI;
\end{lstlisting}

\section{Resources and Collaboration Disclosure}
\begin{itemize}[leftmargin=1.2em]
  \item 02244 Logic for Security lectures and assignment text.
  \item OFMC/AnB tooling and documentation.
  \item AI assistance for writing structure and formatting; final claims validated against OFMC outputs.
\end{itemize}

\end{document}
